\chapter{Arquitetura de Componentes}

\section{Decomposição modular}

O código está dividido em oito módulos JavaScript, cada um com responsabilidade única e
bem delimitada. O \textbf{Component Diagram} é o tipo UML adequado para representar a
divisão de \emph{software} em unidades \emph{deployáveis} e suas dependências de uso.

\diagfiglandscape{diagrams/components.png}{Diagrama de Componentes, com os módulos do Void Descent e suas dependências}{components}

\section{Responsabilidades por módulo}

\begin{description}

  \item[\texttt{src/vd/random.js}] Gerador de números pseudoaleatórios determinístico,
    baseado em \emph{Linear Congruential Generator} (LCG) de 32 bits. É \emph{seedado}
    explicitamente a cada andar (\texttt{random\_seed(0xBADC0DE1 + floor\_id)}), o que
    garante reprodutibilidade do mapa para uma dada combinação de \emph{seed} e andar.

  \item[\texttt{src/vd/sonantx.js}] \emph{Port} JavaScript do sintetizador Sonant.
    Implementa osciladores (sin, quadrada, dente-de-serra, triangular), envelopes ADSR,
    filtros (lopass, hipass, bandpass, notch), LFO e \emph{delay}, todos rodando sobre
    \texttt{Float32Array}.

  \item[\texttt{src/vd/audio.js}] Define os \emph{patches} de instrumentos (oito SFX e a
    faixa ambiente). Inicializa o \texttt{AudioContext}, cria o \texttt{GainNode} mestre
    (\texttt{audio\_master\_gain}) que centraliza o controle de \emph{mute}, e coordena
    a geração assíncrona dos \texttt{AudioBuffer}s no \emph{boot}.

  \item[\texttt{src/vd/renderer.js}] \textit{Pipeline} WebGL. Compila os \emph{shaders}
    de vértice e fragmento, gerencia o atlas de tiles (\texttt{renderer\_bind\_image}),
    expõe primitivas de geometria (\texttt{push\_quad}, \texttt{push\_floor},
    \texttt{push\_block}, \texttt{push\_sprite}) e gerencia até 32 luzes pontuais por
    \emph{frame}.

  \item[\texttt{src/vd/dungeon.js}] Algoritmo de \emph{room placement} com validação
    posterior. Posiciona retângulos não sobrepostos, conecta-os por corredores em forma
    de L, marca tipos especiais (\texttt{start}, \texttt{chest}, \texttt{boss}) e
    converte o resultado em geometria via chamadas a
    \texttt{push\_block}/\texttt{push\_floor}.

  \item[\texttt{src/vd/entities.js}] Hierarquia de classes do mundo simulado:
    classe-base \texttt{entity\_t}, \emph{player}, quatro tipos de inimigos,
    \emph{boss}, projéteis, partículas, explosões e \emph{pickups}. Detalhada no
    Capítulo~\ref{ch:entities}.

  \item[\texttt{src/vd/terminal.js}] \emph{Overlay} HUD textual no estilo terminal de
    boot. Renderiza no elemento DOM \texttt{<code id="a">} (fora do \emph{canvas}) e
    exibe avisos temporários como ``ANDAR 1 :: ANTECAMARA'' ou ``PORTAS TRANCADAS''.

  \item[\texttt{src/vd/game.js}] Módulo orquestrador. Mantém o \texttt{game\_state}
    global, implementa o \texttt{game\_tick} (despachado por estado), e gerencia a
    \emph{shop} de \emph{upgrades}, o \emph{lock-and-clear} de salas, o \emph{spatial
    grid} de colisão e a HUD principal.

\end{description}

\section{Dependências e fluxo de controle}

A arquitetura segue um padrão estritamente \textbf{em camadas}, sem ciclos:

\begin{enumerate}
  \item \textbf{Camada de utilitários puros}, sem dependências:
    \texttt{random.js}, \texttt{sonantx.js}.

  \item \textbf{Camada de infraestrutura}: \texttt{renderer.js}, \texttt{audio.js} (que
    depende de \texttt{sonantx.js}), \texttt{terminal.js}.

  \item \textbf{Camada de geração}: \texttt{dungeon.js}, que depende de
    \texttt{random.js} e da API de geometria de \texttt{renderer.js}.

  \item \textbf{Camada de simulação}: \texttt{entities.js}, que depende das três
    anteriores via globais.

  \item \textbf{Orquestração}: \texttt{game.js}, que depende de tudo.
\end{enumerate}

A ausência de um sistema de módulos (CommonJS, ESM) é intencional. Todos os símbolos
vivem no escopo global definido pela ordem de inclusão dos \texttt{<script>} em
\texttt{index.html}. O padrão minimalista é coerente com a proposta procedural do
projeto e elimina o custo de \emph{boot} de um \emph{bundler}.

\section{Boot sequence}

A inicialização do programa, definida em \texttt{index.html}, segue uma ordem rígida:

\begin{enumerate}
  \item \texttt{audio\_init()} cria o \texttt{AudioContext} e dispara a síntese
    assíncrona da música e dos SFX.
  \item \texttt{renderer\_init()} cria o contexto WebGL, compila os \emph{shaders} e
    aloca o \emph{vertex buffer}.
  \item \texttt{bind\_input()} registra os \emph{handlers} de teclado e mouse no
    \texttt{document}.
  \item \texttt{show\_menu()} define \texttt{game\_state = 'menu'} e instala os
    \emph{handlers} específicos da tela de título.
  \item \texttt{requestAnimationFrame(game\_tick)} inicia o \emph{loop} principal.
\end{enumerate}
